% ---------------------------------------------------------------------------
\section{Introduction}
\label{sec:intro}
% ---------------------------------------------------------------------------

Most real networks are organised into groups of nodes that are more densely connected to
one another than to the rest of the
network~\cite{fortunato2010community,fortunato2022twenty},
and in social, biological and information networks these groups overlap. Methods for
overlapping communities~\cite{xie2013overlapping,fortunato2016community} differ less in
how they search than in what they take a community to be: a group compared with a global
null model~\cite{newman2004finding,traag2019leiden}, with a random
walker's trajectories~\cite{rosvall2008maps}, with a generative
model~\cite{peixoto2014hierarchical,ball2011efficient},
with its significance against a configuration-model null~\cite{lancichinetti2011finding},
or with a consensus among neighbours~\cite{raghavan2007near,gregory2010finding,xie2011slpa}.
No definition is best on every network~\cite{peel2017ground,ghasemian2020evaluating}; the
useful question about a method is which structure it captures, and why.

A natural definition comes from the inside: two friends who share most of their friends
are, more likely than not, in the same circle. This is triadic
closure~\cite{granovetter1973strength}. Real networks are highly
clustered~\cite{watts1998collective}, and neighbourhoods overlap more between members of
a group than across groups, although separating the two effects is itself an inference
problem~\cite{peixoto2022disentangling}. We call the overlap of two closed neighbourhoods
the \emph{commonality} of the pair. Pairwise neighbourhood similarity is an old
tool~\cite{jaccard1912distribution,leicht2006vertex,zhou2009predicting,barzel2013network,feizi2013network},
but a single pair is a weak and degree-confounded signal; what matters is how pairwise
evidence is combined.

We combine it through triples. Three nodes of which one is adjacent to the other two,
and whose three pairs all have commonality at least $t$, form an elementary community;
two such triples belong to the same community when they share a pair. The communities
are the connected components of the two-node-overlap graph of a 3-uniform hypergraph,
which gives the method its name, HyperCommon. Higher-order structures are increasingly
seen as natural units of network
organisation~\cite{benson2016higher,battiston2020networks};
here the hypergraph is built from pairwise commonality rather than given. A node shared
by two communities belongs to triples in both, so overlap needs no extra machinery.

The paper makes three contributions. (i)~\emph{The premise, tested}: on four networks
with recorded memberships, edges inside recorded groups have higher commonality than
edges between them, on the three large ones at every endpoint degree, once unlabelled nodes are no longer
counted as evidence of separation (Sec.~\ref{sec:commonality}). (ii)~\emph{A
construction with exact guarantees}: HyperCommon is independent of processing order,
local and nested in its threshold, every edge it admits belongs to exactly one
community, and on overlapping ring lattices we derive the exact window of thresholds
that recovers the planted communities (Secs.~\ref{sec:method} and~\ref{sec:theory}).
Its cost is linear in the number of wedges, and we show why that number stays near
$m\ln m$ on real networks with hubs (Sec.~\ref{sec:cost}). (iii)~\emph{A map of where it
works}: against fourteen other methods on ring lattices, LFR graphs and seventeen real
networks (Secs.~\ref{sec:protocol} and~\ref{sec:results}), HyperCommon leads where
communities are lattice-like and overlap with unequal sizes and densities, and trails
where communities are random graphs internally, because their members share too few
neighbours.
